\honest{Fake Reductionism}{Eliezer Yudkowsky}{2008}

I quote Keats again: ``We know her woof, her texture.'' I am guessing, ``though it is only a guess,'' that Keats did not know the rainbow's woof and texture himself. Maybe he had only read that Newton explained it as ``light reflected from raindrops,'' which was known in the 13th century. Newton only showed that the light was split into colours rather than changed in colour, but that put rainbows back in the news. \nb{This agrees in substance with the usual history of the rainbow.}

My evidence is a toast. Keats, with Lamb, Wordsworth and Haydon, drank ``Confusion to the memory of Newton'' because ``he destroyed the poetry of the rainbow by reducing it to a prism.'' That is ``one reason to suspect'' that Keats did not understand the subject. \nb{The words come from a letter Haydon wrote to Wordsworth in 1842, nearly twenty-five years after the dinner of December 1817. In Haydon's autobiography the merry Lamb mocks Newton, Keats agrees about the rainbow, and the toast is ``Newton's health, and confusion to mathematics.'' Either way it was a dinner-table joke, which says little about what Keats knew of optics.}

A second guess: Keats could not have explained why the Sun must be behind you, or why the rainbow is an arc of a circle. If so, he had a fake explanation, a ``fake reduction'': he had been told that the rainbow was reduced, but it had not been reduced in his own model of the world.

Anti-reductionists fail to see this distinction between professing that something is reducible and seeing it. They are not much to blame, since it is the general problem I have written about before: knowledge that is only recited, words that stop curiosity, science worn as attire.

Seeing where the rainbow comes from, playing with prisms and making one yourself with a spray of water, is very different from a ``dour-faced philosopher'' telling you that scientists have explained it away, ``Just something to do with raindrops or whatever. Nothing to be excited about.'' I think this difference ``probably accounts for a hell of a lot of the deadly existential emptiness'' said to come with reductionism.

Anti-reductionists experience reduction not as an ``Aha!'' but as being told that the password is ``Science.'' Hearing that science has explained the rainbow only moves it to a genre labelled BORING, by order of the Council of Sophisticated Literary Critics. Their gnomes are pulled out by authority, not dissolved by insight, and nothing beautiful is given in their place, only a sneer that they were fools to find rainbows pretty. ``That's how anti-reductionists experience `reductionism'.'' \nb{The only case given is the guess about Keats. No anti-reductionist or literary critic is quoted describing this experience.}

I excuse Keats: the ``poor lad probably wasn't raised right.'' But he drank to Newton's confusion, so I propose a toast for rationalists: ``To the memory of Keats's confusion.''
